% =====================================================================
% Measuring Equifinality: Conformal Rule Identification in Agent-Based Models
% Submission to CSS 2026 (Santa Fe).
%
% The CSS 2026 call specifies word counts (full paper: 3,000-5,000 words
% inclusive of figures, tables, notes and references) but NO template.
% Accepted full papers appear in Springer Proceedings in Complexity.
%
% This file therefore uses the standard `article' class so that it compiles
% anywhere. To switch to Springer's contributed-volume style when they issue
% it, replace the documentclass line with
%
%     \documentclass{svmult}      % Springer contributed volumes
%  or \documentclass{llncs}       % Springer LNCS
%
% and replace \maketitle/abstract with \title{}\author{}\institute{}
% \maketitle \abstract{...}. Nothing else in the body needs to change.
%
% The bibliography is inline (thebibliography), so a single pdflatex pass
% suffices -- no bibtex/biber run required.
%
% Figures are produced by scripts/fish/13_make_figures.py and live in
% results/figures/fish/ as PDFs. Adjust \graphicspath if you move them.
% =====================================================================

\documentclass[11pt]{article}

\usepackage[margin=1in]{geometry}
\usepackage{graphicx}
\usepackage{booktabs}
\usepackage{amsmath}
\usepackage{amssymb}
\usepackage{microtype}
\usepackage{caption}
\usepackage{siunitx}
% natbib supplies \citep / \citet against the inline thebibliography below,
% whose entries carry explicit [Author(Year)] labels. Load it BEFORE hyperref.
\usepackage[round]{natbib}
\usepackage[hidelinks]{hyperref}

% Figures are generated by scripts/fish/13_make_figures.py. Both paths are
% listed so the file compiles whether the PDFs sit beside the .tex (e.g. after
% uploading to Overleaf) or are read from the repository tree.
\graphicspath{{figures/}{../results/figures/fish/}}

\captionsetup{font=small,labelfont=bf}

% Placeholder for claims pending the verified capacity/depth runs (Phase 1b of
% the capacity reframe). Renders visibly so nothing pending ships silently, and
% uses no extra package.
\newcommand{\todo}[1]{\textbf{[PENDING: #1]}}

\title{\bfseries Measuring Equifinality:\\
Conformal Rule Identification in Agent-Based Models}

\author{[Author names withheld for review]}
\date{}

\begin{document}
\maketitle

% =====================================================================
\begin{abstract}
Agent-based modelers have worried for decades that different behavioral rules can
generate indistinguishable macroscopic patterns. Equifinality is usually stated as a
caveat. We show it can be \emph{measured}.

We recast inverse generative social science as classification over a library of
candidate rules, and attach conformal prediction to the classifier. Conformal
prediction returns a \emph{set} of rules guaranteed to contain the true generator
with probability at least $1-\alpha$, without distributional assumptions and
regardless of how good the classifier is. The size of that set is therefore a
calibrated, observation-conditional measure of equifinality: when the data cannot
separate three mechanisms, the set contains three rules, and the guarantee certifies
that this is a fact about the data rather than a failure of the analyst.

We demonstrate this in three settings. In \textbf{Schelling segregation}, a library
of 2,089 behaviorally distinct utility rules is narrowed to 53 candidates at
$90\%$ confidence -- a 40-fold reduction, with a guarantee. In \textbf{collective
motion}, five agents move under one of eleven known neighbor-selection rules: from
\SI{7}{\second} of trajectory the $95\%$ set contains 3.3 of the eleven rules and
only $1.4\%$ of observations resolve the mechanism uniquely; from \SI{28}{\second} it
contains 1.2 rules and $79\%$ resolve. Equifinality is not a fixed property of a rule
library. It is a property of the library \emph{and the observation}, and it decays
measurably as observation accumulates.

Identifiability is governed by \textbf{what a model can represent}, not simply by
\textbf{how large it is}. Changing how spatial position is \emph{encoded} -- a change
in the model's representational form at a fixed observation and an approximately
matched parameter budget -- inflates the segregation prediction sets by $68\%$. Model
capacity can improve extraction within a model family at a fixed representation, but
cannot distinguish rules that the observation has already rendered identical.
Concretely, a 1.8M-parameter transformer improves on a comparable-budget ($\sim$150k)
one by $+0.142$ at a fixed per-agent representation, yet parameter count alone does not
determine performance -- how the budget is split between width and depth matters too.

Finally we transfer the simulation-trained library onto \textbf{physical robots}
programmed with the same rules. It largely fails ($0.375$), and the reason matters:
robots executing a given rule school $1.7\times$ less tightly than simulated agents
executing it, so the classifier reads their behavior correctly and infers the wrong
rule. Rule identification recovers the \emph{effective} rule of the substrate the
library was built on, not the \emph{programmed} one. Transporting the conformal
calibration across this shift fails silently -- $47\%$ coverage while promising
$95\%$ -- whereas recalibrating on a handful of robot runs restores validity and
returns sets containing 10 of 11 rules: the framework correctly reporting that it
cannot identify the mechanism. \textbf{Conformal prediction guarantees honesty, not
competence.}
\end{abstract}

% =====================================================================
\section{Introduction}

Epstein's generative dictum -- ``if you didn't grow it, you didn't explain it'' --
licenses a model by its ability to produce a target pattern. The inverse question is
harder and more useful: given the pattern, \emph{which rule grew it?} Inverse
generative social science (IGSS) asks exactly this, and typically answers by search
-- evolving candidate rules and keeping those whose output resembles the target
\citep{gunaratne2020,gunaratne2023,epstein2023}.

Two problems dog this program. The first is cost. The second, and deeper, is
\textbf{equifinality}: several distinct rules may generate output that no observer
can tell apart \citep{beven2006}. A search returns a rule; it does not say how many
\emph{other} rules would have served equally well. Modelers know this and hedge, but
the hedge is rhetorical. There is no number.

This paper supplies one. We recast rule recovery as supervised classification over a
library of candidate rules, then wrap the classifier in \textbf{conformal
prediction} \citep{vovk2005,angelopoulos2021}, which converts scores into a set
$C(W)$ of candidate rules with a distribution-free guarantee
%
\begin{equation}
\Pr\bigl(g_{\text{true}} \in C(W)\bigr) \;\ge\; 1-\alpha ,
\label{eq:coverage}
\end{equation}
%
requiring only that calibration and test data be exchangeable. Crucially,
\emph{the guarantee holds whether or not the classifier is any good.} A weak
classifier does not break the coverage; it simply returns larger sets.

That is what makes set size meaningful. Because coverage is guaranteed, $|C(W)|$
cannot be waved away as a modeling artifact: it is a calibrated statement about how
much the \emph{observation} constrains the \emph{rule}. A singleton means the
mechanism is identified. A set of three means three mechanisms remain live at the
requested confidence. \textbf{Equifinality, quantified, conditional on the data you
actually have.}

\paragraph{Contributions.}
\begin{enumerate}
\itemsep2pt
\item \textbf{Conformal sets as an equifinality measure} for ABM rule recovery, with
a coverage guarantee that makes set size interpretable rather than incidental.
\item \textbf{Demonstration across two simulated domains} -- 2,089 rules on a
lattice, 11 rules in continuous space -- with nominal coverage in both, and a
measurement of how equifinality \emph{decays} with observation length.
\item \textbf{An information ladder} read through three parameter layers: which
\emph{observations} make a rule identifiable (relational structure, then temporal
structure), how the model's \emph{representational form} affects resolving power, and
what model \emph{capacity} can and cannot recover -- capacity aids extraction within a
model family at a fixed representation, but cannot distinguish rules the observation
has mapped to identical representations, and the finding that parameter count alone is
not a sufficient description of a model's power.
\item \textbf{A transfer to physical robots} that locates the framework's limit:
rule identification recovers the \emph{effective} rule of its training substrate,
and the conformal layer -- properly recalibrated -- reports that failure honestly
instead of guessing confidently.
\end{enumerate}

% =====================================================================
\section{Background}

\paragraph{Inverse generative social science.}
Evolutionary model discovery \citep{gunaratne2020} evolves rule trees and uses
random-forest importance to rank influential factors; \citet{gunaratne2023} apply it
to residential segregation. The output is a rule and a factor ranking, not a
calibrated set of alternatives.

\paragraph{Equifinality and pattern-oriented modeling.}
That many structures reproduce a pattern is a familiar hazard \citep{beven2006}. The
standard response is pattern-oriented modeling \citep{grimm2005}: constrain the model
with \emph{multiple} patterns, on the reasoning that a rule surviving several filters
is more likely correct. This is sound but informal -- \emph{which} patterns, and how
much constraint does each buy? Section~\ref{sec:ladder} answers that empirically;
Section~\ref{sec:equifinality} turns the residual ambiguity into a number.

\paragraph{Calibration and summary statistics.}
ABM calibration, including approximate Bayesian computation, turns on the choice of
summary statistics: statistics that discard the information distinguishing two
mechanisms render those mechanisms inferentially identical. Our information ladder
measures this directly.

\paragraph{Conformal prediction.}
Split conformal \citep{vovk2005} uses a held-out calibration set to convert scores
into sets with finite-sample coverage. RAPS \citep{angelopoulos2021} is an adaptive
variant penalizing long tails. We use RAPS throughout.

% =====================================================================
\section{Framework}

The pipeline is domain-agnostic and has four stages.

\begin{enumerate}
\itemsep2pt
\item \textbf{Rule library.} Enumerate or generate a set $\mathcal{G}$ of candidate
behavioral rules, and deduplicate them \emph{behaviorally} -- two rules that produce
indistinguishable dynamics are one rule, and keeping both would manufacture ambiguity
rather than measure it.
\item \textbf{Forward simulation.} Run each rule to produce labelled observations.
\item \textbf{Recovery.} Train a classifier $P(g \mid W)$ mapping an observation
window $W$ to a distribution over rules.
\item \textbf{Conformal layer.} On a held-out calibration partition never used for
training or model selection, calibrate RAPS to produce $C(W)$ satisfying
Equation~\ref{eq:coverage}.
\end{enumerate}

\paragraph{Three kinds of parameter, and the pipeline as a map.}
Rule-recovery studies routinely conflate three distinct kinds of parameter, and
separating them is what lets us ask \emph{where} identifiability is won or lost.
\emph{Generative-system} parameters $\theta_{\mathrm{ABM}}$ define the world being
simulated -- interaction strengths, neighborhood radius, density, agent count, and the
update rule itself, which here is the target of inference. \emph{Observation and
representation} parameters $\theta_{\mathrm{obs}}$ determine how that world becomes
evidence -- observation length, sampling frequency, aggregation functions, token
definitions, feature invariances, and whether individual-agent identity is retained.
\emph{Inference-model} parameters $\theta_{\mathrm{ML}}$ define the learner -- depth,
width, attention heads, recurrent-state size, total trainable parameters, and the
model's internal representational form. The four stages above then compose into a map
%
\begin{equation}
\theta_{\mathrm{ABM}} \;\longrightarrow\; X
  \;\xrightarrow{\;\phi_{\mathrm{obs}}\;} Z
  \;\xrightarrow{\;f_{\theta_{\mathrm{ML}}}\;} \widehat{\Gamma}_{\alpha}(Z),
\label{eq:pipeline}
\end{equation}
%
where $X$ is the simulated trajectory, $\phi_{\mathrm{obs}}$ the
observation/tokenization map, $Z$ the machine-readable representation, $f$ the
classifier, and $\widehat{\Gamma}_{\alpha}(Z)$ the conformal rule set $C(W)$ of
Equation~\ref{eq:coverage}. The distinction is not pedantic: a recurring error is to
treat a larger $\theta_{\mathrm{ML}}$ as if it could recover information that
$\phi_{\mathrm{obs}}$ discarded upstream. It cannot. \emph{Capacity cannot distinguish
states or rules that the observation has mapped to identical representations}; what a
model can recover is bounded by $\theta_{\mathrm{obs}}$ first, and by the form -- not
merely the size -- of $\theta_{\mathrm{ML}}$ second. (For Schelling, $\theta_{\mathrm{ABM}}$
is under our control in-repo; for collective motion we use released simulated
trajectories with known rules, so $\theta_{\mathrm{ABM}}$ is a property of the supplied
data rather than a simulator we vary.)

\paragraph{The experimental unit is the simulation run, never the observation
window.} Windows from the same run are dependent; treating them as independent
inflates apparent performance and silently voids the conformal guarantee, which needs
calibration and test scores to be exchangeable. Throughout, whole runs are assigned
to partitions, evaluation windows within a run do not overlap, and all confidence
intervals bootstrap over runs.

% =====================================================================
\section{Case study 1: Schelling segregation (2,089 rules)}
\label{sec:schelling}

\paragraph{Rules.}
In Schelling's model \citep{schelling1971} an agent relocates when dissatisfied with
its neighborhood; what makes it a rule-recovery problem is that the utility function
driving that decision is not observed. Following \citet{gunaratne2023}, we take that
utility to be an arithmetic expression tree over six primitive factors -- same-type fraction, neighbor tenure,
distance from home, vacancy fraction, recent moves, and neighbor utility -- combined
with $+\,-\,\times\,\div$. From 10,000 generated candidates (depth $\le 3$),
behavioral deduplication against fixed probe simulations retains \textbf{2,089
behaviorally distinct trees}, which become the class labels.

\paragraph{Observations.}
A $50\times50$ torus, density $0.95$, Moore neighborhood, 500 ticks, with snapshots
at ticks 100--500. Each snapshot is tokenized into $2\times2$ patches; a transformer
encoder (3.9M parameters) with decomposed time/space embeddings classifies the
sequence. Each tree is simulated 100 times, giving ${\sim}209$k labelled sequences
split 80/10/10 by run.

\paragraph{Results.}
Chance top-1 is $1/2089 = 0.048\%$. The classifier reaches top-1 $0.319$ ($665\times$
chance), top-3 $0.485$, top-5 $0.562$, with pre-temperature ECE $0.017$: it extracts
a strong signal about the generating rule from five grid snapshots, and its raw
softmax is already well calibrated.

\begin{table}[t]
\centering
\small
\caption{Conformal sets on the Schelling rule library. \textbf{At $90\%$ confidence,
2,089 candidate rules are narrowed to 53: a 40-fold reduction, with a guarantee.}}
\label{tab:schelling-conformal}
\begin{tabular}{lccc}
\toprule
$\alpha$ & target & empirical coverage & mean $|C|$ (\% of library) \\
\midrule
0.01 & 0.99 & 0.989 & 234.3 \; (11.2\%) \\
0.05 & 0.95 & 0.948 & 104.2 \; (5.0\%) \\
0.10 & 0.90 & 0.897 & \textbf{52.9} \; (\textbf{2.5\%}) \\
0.20 & 0.80 & 0.801 & 20.5 \; (1.0\%) \\
\bottomrule
\end{tabular}
\end{table}

Coverage tracks the nominal level at every $\alpha$ (Table~\ref{tab:schelling-conformal}).
In absolute terms the sets are large -- but the right denominator is the library. This
is a quantitative statement about how much a five-snapshot observation of a
segregation process constrains the underlying utility rule, and it is the kind of
statement the IGSS literature has not previously been able to make. The remaining
ambiguity is not noise but structure: the surviving rules are behaviorally similar
trees.

\subsection{Representational form matters, at fixed observation}
\label{sec:rowcol}

We ablate one thing: how spatial position is encoded. The base model learns a
separate embedding for each of the 625 patch positions; the ablation decomposes it
into a learned row vector plus a learned column vector. This \emph{reduces} parameters
slightly (3.73M vs 3.88M) and leaves the observation, the data, and the training
budget untouched: the representation $Z$ the model consumes is bit-for-bit identical.
Only the model's internal representational form -- a choice inside
$\theta_{\mathrm{ML}}$, not a change to $\theta_{\mathrm{obs}}$ -- differs.

\begin{table}[t]
\centering
\small
\caption{Spatial-embedding ablation. Accuracy barely moves; the conformal sets
inflate by $68\%$. A model can be equally \emph{right} and considerably less
\emph{informative}.}
\label{tab:rowcol}
\begin{tabular}{lccc}
\toprule
 & top-1 & mean $|C|$ @ $\alpha{=}0.10$ & @ $\alpha{=}0.05$ \\
\midrule
flat spatial embedding      & \textbf{0.319} & \textbf{52.9} & \textbf{104.2} \\
decomposed row $+$ column   & 0.312 & 88.7 & 192.4 \\
\bottomrule
\end{tabular}
\end{table}

Accuracy is nearly unchanged ($0.319$ vs $0.312$), but the row/column model needs 89
candidate rules where the flat model needs 53 to make the same coverage guarantee.
This matters twice. First, \textbf{set size is a more discriminating diagnostic than
accuracy} -- a difference invisible in top-1 is stark in the equifinality measure.
Second, it isolates a distinct lever: this is a change in the model's
\emph{representational form} (a choice within $\theta_{\mathrm{ML}}$) at a fixed
observation $Z$ and an approximately matched parameter budget, and it costs $68\%$ of
the resolving power. This is a different axis from model capacity: in
Section~\ref{sec:ladder}, capacity within a model family (at a fixed per-agent
representation) moves accuracy by up to $+0.142$; here, at a fixed observation and
near-matched budget, only the representational form changes, and it costs $68\%$ of the
resolving power.

% =====================================================================
\section{Case study 2: collective motion (11 rules)}
\label{sec:fish}

The segregation study leaves an objection: perhaps the classifier recognizes
\emph{grid textures}, not behavioral rules. So we move to a domain with no grid, no
lattice, and a rule space small enough to be exhaustively understood.

\subsection{The benchmark}

Five agents move in a circular arena ($R = \SI{0.25}{\metre}$) under burst-and-coast
dynamics. A rule $g=(s,k)$ fixes a neighbor-selection strategy $s$ and the number $k$
of neighbors attended to: no interaction ($k{=}0$); the $k$ \textbf{nearest}, $k$
\textbf{random}, or $k$ \textbf{most influential} neighbors for $k \in \{1,2,3\}$; or
\textbf{all} four. Eleven rules; chance $0.091$.

The families differ \emph{only} in which neighbors enter the social sum -- the
interaction functions are identical. This makes the benchmark sharp, and it makes the
high-$k$ rules genuinely hard: with $k=3$ of 4 available neighbors, ``nearest three'',
``random three'' and ``most influential three'' draw from nearly the same pool.

We use the \textbf{simulated-agent} trajectories released with \citet{lei2020}, who
build on the burst-and-coast model of \citet{calovi2018}. These are model output, not
empirical fish data, and the point is precisely that: the generating rule is known for
every run. Eleven rules $\times$ 50 runs $=$ 550 runs, balanced. Time is irregular
\textbf{kick time} (median interval \SI{0.434}{\second}), so we index sequences by
kick event rather than resampled clock time -- interpolating onto a uniform grid would
smooth away the discontinuities through which the rule acts.

\subsection{A leakage hazard worth documenting}

Experiment identifiers are \textbf{shared across the eleven rule files}, and runs
begin from common initial placements: inspecting the first recorded event of every
run, one experiment is identical across all eleven rules and thirteen more are
partially identical. A classifier could learn the \emph{initial condition} rather than
the rule.

We block this twice: the grouping key is the experiment id applied across all eleven
rules at once (if experiment 42 is in test, all eleven of its runs are), and we
discard a 50-event burn-in that removes both the placement transient and the region
where the coupling lives. Grouped 5-fold cross-validation over the 50 seeds; each
experiment is tested exactly once. We report this because it is the kind of hazard
that inflates results silently, and because it is a property of a widely-used public
dataset.

\subsection{The information ladder}
\label{sec:ladder}

We compare four \textbf{observation designs}, holding the task fixed, at $T=64$ kicks
(\SI{28}{\second}); see Figure~\ref{fig:ladder} and Table~\ref{tab:ladder}.

\begin{table}[t]
\centering
\small
\caption{The information ladder: two large observation gains. The final row compares
a different model family at higher capacity; see text and Section~\ref{sec:ladder} for
what that comparison does and does not license.}
\label{tab:ladder}
\begin{tabular}{llc}
\toprule
What is observed & Model & Exact-rule accuracy \\
\midrule
nothing                    & chance                      & 0.091 \\
solo kinematics            & random forest               & $0.352 \pm 0.008$ \\
$+$ relational             & random forest               & $0.717 \pm 0.012$ \\
$+$ temporal               & GRU (156k params)           & $\mathbf{0.929 \pm 0.023}$ \\
$+$ per-agent granularity  & transformer (1.8M params)   & $0.927 \pm 0.020$ \\
\bottomrule
\end{tabular}
\end{table}

\begin{figure}[t]
\centering
\includegraphics[width=0.86\textwidth]{fig1_ladder.pdf}
\caption{Identifiability is set largely by what the model observes. Relational
information is worth $+0.365$; the time course a further $+0.212$. Model capacity is
examined separately (Table~\ref{tab:capacity}): it helps within a model family but does
not substitute for what the observation retains.}
\label{fig:ladder}
\end{figure}

\textbf{Relational information is worth $+0.365$.} An observer recording only how
each fish moves -- speed, turning, wall distance -- recovers the rule $35\%$ of the
time. Letting it see the neighbors doubles that.

\textbf{Temporal information is worth $+0.212$.} The \emph{identical} relational
statistics, fed to a model that sees their time course rather than their window
average, go from $0.717$ to $0.929$. The rules differ in the \emph{persistence} of the
influencing set, and a window average cannot express persistence.

\begin{table}[t]
\centering
\small
\caption{Model capacity at a \emph{fixed} per-agent representation ($T=64$). At a fixed
four-layer depth, raising the parameter budget buys $+0.142$; a comparable-budget depth
increase does not. The GRU, handed engineered group statistics, answers a different
(representational-burden) question. Exact-rule accuracy; confidence intervals
exp\_id-clustered over the 50 experiments.}
\label{tab:capacity}
\begin{tabular}{llc}
\toprule
Model & Parameters & Exact-rule accuracy \\
\midrule
transformer, 2 layers, $d{=}48$  & 148{,}091     & 0.803 \\
transformer, 4 layers, $d{=}40$  & 182{,}771     & 0.785 \\
transformer, 4 layers, $d{=}128$ & 1{,}812{,}619 & \textbf{0.927} \\
GRU (group statistics)           & 155{,}659     & \textbf{0.929} \\
\bottomrule
\end{tabular}
\end{table}

\textbf{Capacity matters within a model family, but parameter count is not a
sufficient description.} The three transformers all receive the \emph{same} per-agent
representation, so comparing them isolates what model capacity recovers
(Table~\ref{tab:capacity}). Holding the per-agent representation and the four-layer
architecture fixed, increasing the parameter budget from 182{,}771 to 1{,}812{,}619
raised test accuracy from $0.785$ to $0.927$ -- a paired improvement of $+0.142$
($95\%$ CI $[+0.129,+0.156]$, exp\_id-clustered; run-clustered $[+0.132,+0.154]$).
Capacity is therefore not free, and the earlier reading that an ``eleven-fold increase
buys nothing'' conflated this within-family effect with a cross-architecture tie. But
parameter count alone does not determine performance either: at a comparable budget, a
four-layer $d{=}40$ model ($0.785$) did not improve on -- and was slightly below -- a
two-layer $d{=}48$ model ($0.803$; paired $-0.018$, $95\%$ CI $[-0.028,-0.008]$). Model
size cannot be read independently of how the budget is allocated across width and
depth, nor of the observation and tokenization upstream.

\textbf{The GRU tie is a difference of representational burden.} The 156k-parameter
GRU matches the 1.8M-parameter transformer ($0.929$ vs $0.927$; paired $+0.0015$,
$95\%$ CI $[-0.003,+0.006]$) not because capacity is irrelevant but because the two
carry different burdens: the GRU is handed thirteen engineered group statistics, while
the transformer must discover comparable aggregates from raw per-agent tokens. They
share overlapping task-relevant information at different representational burden -- a
statement about representation, not a contest between architectures.

\paragraph{Training to convergence is a fairness prerequisite, not one of the factors
under study.} Our first transformer, trained for 40 epochs, scored $0.825$ and
appeared to lose to the GRU by ten points. Its validation accuracy was still climbing
when the epoch budget ran out; early stopping never fired. At 120 epochs the same
architecture reaches $0.927$. \textbf{A comparison across model families is worthless
unless every model is trained to convergence} -- convergence is a fairness and
reproducibility requirement, held separate from the structural factors (observation,
representation, capacity) the comparison is about. The failure mode is silent: an
undertrained model returns a perfectly plausible number, and a comparison against it
will confidently support the wrong conclusion. The two capacity control runs above were
trained under this same 120-epoch protocol and reached convergence -- validation
plateaued over the final 30 epochs, with early stopping firing in two folds of the
$d{=}40$ model -- so Table~\ref{tab:capacity} compares converged models.

\subsection{Identifiability grows with observation, then saturates}

\begin{table}[t]
\centering
\small
\caption{Exact-rule accuracy against observation length. Every rung rises, and the
gaps between rungs persist.}
\label{tab:window}
\begin{tabular}{llccc}
\toprule
$T$ (kicks) & seconds & solo kinematics & $+$ relational (RF) & $+$ temporal (GRU) \\
\midrule
16  & 7   & 0.254 & 0.549 & \textbf{0.771} \\
32  & 14  & 0.291 & 0.624 & 0.869 \\
64  & 28  & 0.352 & 0.717 & 0.929 \\
128 & 56  & 0.434 & 0.809 & 0.967 \\
256 & 111 & 0.533 & 0.885 & \textbf{0.984} \\
512 & 222 & 0.661 & 0.962 & 0.970 \\
\bottomrule
\end{tabular}
\end{table}

With two minutes of trajectory the sequence model is at $0.984$ and the task is
effectively solved (Table~\ref{tab:window}). Rule identification here is therefore
\emph{easy given enough observation}, which means the scientifically interesting
regime is the \textbf{short} one, where the mechanism is genuinely underdetermined.
That is also where conformal prediction earns its keep.

\subsection{Equifinality decays as observation accumulates}
\label{sec:equifinality}

This is the central measurement of the paper (Figure~\ref{fig:equifinality},
Table~\ref{tab:equifinality}).

\begin{table}[t]
\centering
\small
\caption{RAPS at $95\%$ confidence, across observation length. Coverage is
\emph{tight} ($0.949$ against a $0.950$ target) at $T=16$, where the model is
genuinely uncertain and the guarantee is actually tested.}
\label{tab:equifinality}
\begin{tabular}{llccc}
\toprule
$T$ & seconds & accuracy & empirical coverage & mean $|C|$ \; (singletons) \\
\midrule
\textbf{16} & 7   & 0.771 & \textbf{0.949} \; [0.944, 0.955] & \textbf{3.28} \; (1.4\%) \\
32          & 14  & 0.869 & 0.957 & 1.86 \; (44\%) \\
64          & 28  & 0.929 & 0.968 & 1.22 \; (79\%) \\
128         & 56  & 0.967 & 0.983 & 1.07 \; (94\%) \\
256         & 111 & 0.984 & 0.989 & 1.02 \; (98\%) \\
512         & 222 & 0.970 & 0.986 & 1.06 \; (94\%) \\
\bottomrule
\end{tabular}
\end{table}

\begin{figure}[t]
\centering
\includegraphics[width=\textwidth]{fig6_equifinality.pdf}
\caption{\textbf{Equifinality is measurable, and it decays with observation.} Left:
mean prediction-set size against observation length. Right: how many rules remain
live at $90\%$ confidence. At seven seconds, two-thirds of observations cannot pin
the mechanism to a single rule.}
\label{fig:equifinality}
\end{figure}

\textbf{At seven seconds of observation, the average $95\%$-confidence prediction set
contains 3.28 of the eleven rules, and only $1.4\%$ of observations resolve the
mechanism to a single rule. By twenty-eight seconds the set holds 1.22 rules and
$79\%$ resolve.} Equifinality is therefore \textbf{not a fixed property of a rule
library}. It is a property of the library \emph{and the observation}, and it decays
measurably as observation accumulates. That is the quantity this paper proposes to
report.

\paragraph{The over-coverage at long windows is a ceiling effect, and it is
instructive.} With top-1 accuracy at $0.93$ and a convention of never returning an
empty set, coverage cannot fall below ${\sim}0.93$ whatever $\alpha$ requests; the
sets have collapsed to singletons and there is nothing left to express.
\emph{Measuring equifinality at an observation length where the task is already solved
measures nothing.} This is a caution for anyone tempted to calibrate on a saturated
model and conclude that a system exhibits no equifinality: they will have measured the
ceiling, not the system.

\subsection{The residual equifinality is structured}

We expected confusions to stay inside a strategy family and slip on $k$
(``nearest-2'' mistaken for ``nearest-3''). \textbf{The reverse is true.} The
transformer recovers the exact rule at $0.927$, the strategy family at $0.939$, and
the \emph{neighbor count} $k$ at $\mathbf{0.959}$. Of its errors, $43\%$ preserve the
correct $k$ while getting the family wrong; only $16\%$ do the opposite.

\textbf{How many neighbors an agent attends to is far easier to recover than which
neighbors it selects.} The mechanism explains it. $k$ sets the effective strength of
cohesion, written directly into group geometry -- mean nearest-neighbor distance falls
monotonically with $k$ within every family ($0.28$, $0.22$, $0.19$ for
nearest-1/2/3). \emph{Which} neighbors are chosen is a subtler, higher-order property,
and at $k=3$ of four available neighbors the three families draw from an almost
identical pool. They converge behaviorally. No observer -- and no amount of model
capacity -- can separate them.

\begin{figure}[t]
\centering
\includegraphics[width=0.62\textwidth]{fig4_confusion.pdf}
\caption{The errors are not scattered. They concentrate in a high-$k$ cluster --
nearest-3, random-3, influential-3 and all-4 -- where the selection rules draw from
nearly the same pool of neighbors and therefore converge behaviorally. No-interaction
is recovered perfectly.}
\label{fig:confusion}
\end{figure}

The confusion matrix (Figure~\ref{fig:confusion}) shows this directly: the errors
concentrate in the high-$k$ cluster, while no-interaction is recovered perfectly and
random-1 at $0.97$.

\textbf{This is where the conformal set pays off.} A prediction set
$C(W) = \{\text{nearest-3},\ \text{influential-3},\ \text{all-4}\}$ is not a
classifier failing. It is a calibrated, coverage-guaranteed statement that
\emph{three mechanisms are behaviorally indistinguishable given this observation}. The
equifinality is a property of the system and the data, not of the analyst's model
choice -- and the guarantee is what licenses that reading. Without it, a large
prediction set would be indistinguishable from a weak classifier.

% =====================================================================
\section{Case study 3: transfer to physical robots}
\label{sec:robots}

Everything so far is simulation. The framework's real claim is that a rule library
learned in simulation identifies mechanisms in a system that was \emph{not} simulated.
The same release lets us test exactly that: five \textbf{cuboid robots}, each
programmed with one of the same neighbor-selection strategies, moving in a larger
arena. Real hardware, real actuation error, real collisions -- but the generating rule
is still known, because it was programmed.

The domain shift is substantial: the arena is $1.7\times$ larger
(\SI{0.42}{\metre} vs \SI{0.25}{\metre}), kicks are $1.9\times$ slower
(\SI{0.840}{\second} vs \SI{0.434}{\second}), runs are shorter, and only ten of the
eleven rules have robot files. We take the transformer already trained on simulated
agents, \textbf{change nothing about it}, and run it on 1,961 non-overlapping robot
windows.

\paragraph{Units first.}
Positions are normalized by arena radius, so lengths are already dimensionless -- but
time is not. Speed is in arena-radii per \emph{second}, and robots kick $1.9\times$
more slowly, so a robot moving identically (the same fraction of the arena per kick)
registers about half the speed. We therefore nondimensionalize, adopting each domain's
own kick interval $\tau$ as its time unit: $\text{speed}\times\tau$,
$\text{accel}\times\tau^{2}$, $\mathrm{d}t/\tau$. This is a unit conversion, not
adaptation; it retrains nothing. It is worth $+0.085$ top-3 accuracy, and skipping it
would be as wrong as adding feet to metres.

\subsection{The transfer largely fails -- and the reason is instructive}

The same model scores $0.927$ on simulated agents and $\mathbf{0.375}$ on robots
(chance $0.10$). It retains real signal ($3.7\times$ chance) but is far from its
simulation performance. The \emph{pattern} of failure is the clue: per-rule recall is
$0.89$ for no-interaction and $0.74$ for nearest-1, and collapses to ${\sim}0$ for
every rule with $k \ge 2$, reaching exactly $0$ for all-neighbors. \textbf{The model
reads the robots as weakly interacting.}

It is right to. Comparing mean nearest-neighbor distance for the \emph{same rule}
across substrates (Table~\ref{tab:cohesion}), the discrepancy is \textbf{zero when
there is no interaction} and \textbf{grows monotonically with how cohesive the rule
ought to be}. Physical embodiment -- inertia, actuation error, sensing delay,
collision avoidance -- acts as a low-pass filter on interaction strength. A robot
programmed with all-neighbors achieves the cohesion of a simulated agent running
nearest-1.

\begin{table}[t]
\centering
\small
\caption{Mean nearest-neighbor distance (arena radii) for the \emph{same rule} on
each substrate. The gap is zero without interaction and grows with the rule's
cohesiveness.}
\label{tab:cohesion}
\begin{tabular}{lccc}
\toprule
rule & agents & robots & ratio \\
\midrule
no-interaction & 0.549 & 0.509 & \textbf{0.93$\times$} \\
nearest-1      & 0.280 & 0.406 & 1.45$\times$ \\
nearest-2      & 0.215 & 0.338 & 1.57$\times$ \\
influential-2  & 0.209 & 0.307 & 1.47$\times$ \\
all-4          & 0.182 & 0.313 & \textbf{1.71$\times$} \\
\bottomrule
\end{tabular}
\end{table}

So the classifier is not malfunctioning. It is correctly reading the \emph{behavior},
and the behavior says ``weak interaction''. It is right about the dynamics and wrong
about the rule, because \textbf{the same rule produces different dynamics on different
hardware}.

This is the paper's sharpest caution for IGSS. Behavioral-rule identification recovers
the \textbf{effective} rule -- the rule that would produce this behavior in the
substrate the library was built on -- not the \textbf{programmed} rule. A rule library
cannot be transported across embodiments without accounting for the substrate's own
dynamics.

\subsection{The conformal layer reports the failure honestly}

\begin{table}[t]
\centering
\small
\caption{Conformal prediction under domain shift. Transporting the calibration set
across the shift fails \emph{silently}: the sets get \emph{smaller}, not larger.
Recalibrating on held-out robot runs restores validity with no retraining.}
\label{tab:robot-conformal}
\begin{tabular}{llcccc}
\toprule
 & & \multicolumn{2}{c}{calibrated on \textbf{simulation}}
   & \multicolumn{2}{c}{calibrated on \textbf{robots}} \\
\cmidrule(lr){3-4}\cmidrule(lr){5-6}
$\alpha$ & target & coverage & mean $|C|$ & coverage & mean $|C|$ \\
\midrule
0.05 & 0.95 & \textbf{0.473} & 1.6 & \textbf{0.939} & 10.0 \\
0.10 & 0.90 & 0.440 & 1.4 & 0.898 & 9.3 \\
0.20 & 0.80 & 0.405 & 1.2 & 0.815 & 6.9 \\
\bottomrule
\end{tabular}
\end{table}

\textbf{Transporting the calibration set across the domain shift fails silently and
catastrophically} (Table~\ref{tab:robot-conformal}). It returns small, confident sets
of 1.2--1.6 rules that contain the truth $47\%$ of the time while promising $95\%$.
Nothing in the output warns you. Exchangeability -- the one assumption conformal
prediction actually needs -- has been violated, and the guarantee is void. This is the
trap, and it is easy to fall into: the numbers look \emph{better}, not worse.

\textbf{Recalibrating on held-out robot runs restores validity} -- same weights, no
retraining, only the conformal quantile re-estimated from runs of the new system.
Coverage returns to nominal. And the restored sets contain \textbf{10 of the 11
rules}.

That is not a failure of the method. It is the method working. The set size is a
calibrated statement, and the statement it makes is: \emph{given a library built in
simulation, these robot trajectories do not identify the mechanism.} The framework
declines to guess rather than guessing confidently and wrongly.

\textbf{Conformal prediction guarantees honesty, not competence.} It cannot
manufacture information that the observation does not contain, and -- properly
calibrated -- it does not pretend to.

% =====================================================================
\section{Discussion}

\paragraph{Equifinality becomes an estimand.}
Across all three settings, the conformal set converts a qualitative caveat into a
measured quantity with an interpretable unit: the number of rules surviving at
confidence $1-\alpha$, computable per observation. The measure behaves sensibly at
very different scales -- narrowing 2,089 segregation rules to 53, eleven movement
rules to one or two, and (correctly) declining to narrow the robot rules at all. One
can now ask, quantitatively: \emph{how much data must I collect before mechanism A and
mechanism B become distinguishable?}

\paragraph{Observation and representation shape identifiability; capacity aids
extraction within a model family.} Read through the three-parameter split of
Equation~\ref{eq:pipeline}, the case studies locate where identifiability is won and
lost. Segregation varies the model's \emph{representational form} ($\theta_{\mathrm{ML}}$)
at a fixed observation $Z$ and an approximately matched budget: accuracy is nearly
identical, yet the conformal sets differ by $68\%$. Collective motion varies model
\emph{capacity} ($\theta_{\mathrm{ML}}$ size) at a fixed per-agent representation: at a
fixed four-layer depth, raising the parameter budget from 182{,}771 to 1{,}812{,}619
improves accuracy by $+0.142$, but a comparable-budget \emph{depth} increase does not
(Table~\ref{tab:capacity}). The unifying statement is that capacity can improve
extraction within a model family when the representation is fixed, but cannot
distinguish rules that the observation ($\theta_{\mathrm{obs}}$) has already mapped to
identical representations -- and that model size is not a sufficient description of a
model's power, since parameter count cannot be read independently of how it is
allocated across width and depth, of tokenization, or of the observation upstream. In
practice this means an IGSS study should report not only a parameter count but the
tokenization, observation structure, width and depth, and whether scientific
invariances are engineered or learned.

We would not have learned this without the aggregate control. A transformer beating a
random forest looks like a victory for agent-level modeling; the same transformer
beside a GRU -- which is handed the group statistics the transformer must discover from
per-agent tokens -- reveals the gain came from the time axis, and reframes the
GRU/transformer contrast as one of differing representational \emph{burden} over
overlapping task-relevant information, not a contest between architectures.

\paragraph{A cautionary note on conformal implementations.}
We found a defect in a common RAPS formulation: if the randomization term in the
non-conformity score is applied only when the true class is \emph{not} ranked first,
then for an accurate model most calibration points receive an inflated score, the
calibration quantile inflates with them, and every prediction set comes out too large.
Coverage remains \emph{valid}, so nothing in the output announces the problem. Notably,
\textbf{the severity scales with model accuracy}: our collective-motion model ($93\%$
top-1) had ${\sim}93\%$ of calibration points at rank 0 and its coverage sat at $0.97$
regardless of $\alpha$; our segregation model ($32\%$ top-1) had only ${\sim}32\%$, and
its numbers barely moved. The bug is most dangerous exactly where the model is best --
which is when one is least inclined to be suspicious. For any work measuring
equifinality \emph{by set size}, this is a silent and material error.

\paragraph{Reproducibility.}
Every number reported here rests on a five-fold grouped cross-validation in which each
experiment is tested exactly once, with all confidence intervals bootstrapped over
runs rather than windows. As an independent check, the sequence model was retrained
from scratch on separate hardware (an NVIDIA H100 on a cluster, against an RTX 4090
workstation): at $T=128$ the two runs give $0.9640$ and $0.9671$ -- the same number
within fold-to-fold noise.

\paragraph{Limitations.}
The two rule-recovery case studies use simulated data with known ground truth; the
robot study uses real hardware but a programmed, not discovered, rule. The Schelling
library, though large, is a depth-3 subset of an infinite rule space. Exchangeability
holds at the level of runs rather than windows; we mitigate structurally and bootstrap
over runs. With 50 seeds, the fold-to-fold spread in the collective-motion study is
about $\pm 0.02$, which is the honest resolution of those numbers. Finally, the
transformer is trained to convergence only at $T=64$; the observation-length sweep
uses the sequence model, which the paired comparison in Section~\ref{sec:ladder}
licenses, and we make no claim about the transformer at other window lengths.

\paragraph{Next.}
The same release contains trajectories of \emph{real fish}. Section~\ref{sec:robots}
suggests the obvious hazard -- a library built in simulation will recover the effective
rule of the simulator, not of the animal -- and equally the obvious remedy: calibrate on
the target substrate, and let the conformal set report how much, or how little, the
observation actually determines.

% =====================================================================
\begin{thebibliography}{9}
\small

\bibitem[Angelopoulos et al.(2021)]{angelopoulos2021}
Angelopoulos, A., Bates, S., Malik, J., Jordan, M. (2021).
Uncertainty sets for image classifiers using conformal prediction.
\emph{ICLR}.

\bibitem[Beven(2006)]{beven2006}
Beven, K. (2006).
A manifesto for the equifinality thesis.
\emph{Journal of Hydrology}, 320(1--2), 18--36.

\bibitem[Calovi et al.(2018)]{calovi2018}
Calovi, D.\,S., Litchinko, A., Lecheval, V., Lopez, U., P\'erez Escudero, A.,
Chat\'e, H., Sire, C., Theraulaz, G. (2018).
Disentangling and modeling interactions in fish with burst-and-coast swimming.
\emph{PLoS Computational Biology}, 14(1), e1005933.

\bibitem[Epstein(2023)]{epstein2023}
Epstein, J.\,M. (2023).
Inverse generative social science: backward to the future.
\emph{JASSS}, 26(2).

\bibitem[Grimm et al.(2005)]{grimm2005}
Grimm, V., Revilla, E., Berger, U., Jeltsch, F., Mooij, W.\,M., Railsback, S.\,F.,
Thulke, H.-H., Weiner, J., Wiegand, T., DeAngelis, D.\,L. (2005).
Pattern-oriented modeling of agent-based complex systems.
\emph{Science}, 310(5750), 987--991.

\bibitem[Gunaratne and Garibay(2020)]{gunaratne2020}
Gunaratne, C., Garibay, I. (2020).
Evolutionary model discovery of causal factors behind the socio-agricultural
behavior of the Ancestral Pueblo.
\emph{PLOS ONE}.

\bibitem[Gunaratne et al.(2023)]{gunaratne2023}
Gunaratne, C., Rand, W., Garibay, I. (2023).
Generating mixed patterns of residential segregation: an evolutionary approach.
\emph{JASSS}, 26(2).

\bibitem[Lei et al.(2020)]{lei2020}
Lei, L., Escobedo, R., Sire, C., Theraulaz, G. (2020).
Computational and robotic modeling reveal parsimonious combinations of interactions
between individuals in schooling fish.
\emph{PLoS Computational Biology}, 16(3), e1007194.

\bibitem[Schelling(1971)]{schelling1971}
Schelling, T.\,C. (1971).
Dynamic models of segregation.
\emph{Journal of Mathematical Sociology}, 1(2), 143--186.

\bibitem[Vovk et al.(2005)]{vovk2005}
Vovk, V., Gammerman, A., Shafer, G. (2005).
\emph{Algorithmic Learning in a Random World}. Springer.

\end{thebibliography}

\end{document}
